\begin{center}
\large\textsf{\textbf{Electronic Engineering}}\vspace{0.5ex}\hrule\vspace{0.5ex}
\large\textsf{Department of Electronic Engineering}
\end{center}
% ----------------------------------------------------- %
% DEGREE PROGRAMME INFO
\subsection*{Degree Programme Information}
\vspace{0ex}%
\noindent\parbox{0.5\textwidth}{%
    \noindent\begin{tabular}{@{}ll}
                 \textsf{Programme:}          & 	Electronic Engineering  \\
                 \textsf{Programme Code:}         & 	TEE    \\
                 \textsf{Duration:}         & 	5 years    \\
                 \textsf{Minimum Credits:}    & 	600 \\
                 \textsf{Maximum Credits:}    & 	675 \\
                 \textsf{ZNQF Level:}    & 	SADC-QF-Level 8
    \end{tabular}}
\vspace{2ex}\hrule\vspace{2ex}


% ----------------------------------------------------- %
% ENTRY REQUIREMENTs
\subsection*{Entry Requirements}
\begin{outline}
    \1      \textbf{Normal entry} \newline Requires at three ‘A’ Level passes in Pure/Additional/Mechanical Mathematics, Physics and Chemistry.
    \1      \textbf{Special entry requires:}
    \begin{enumerate}[label=\alph*)]
        \item National Diploma in Electronic Engineering or Telecommunication Engineering or Instrumentation and Control or Computer Engineering plus 2 years post National Diploma working experience.
        \item Higher National Diploma in Electronic Engineering or Telecommunication Engineering or Instrumentation and Control or Computer Engineering plus 1 year post Higher National Diploma working experience.
    \end{enumerate}
\end{outline}
\vspace{2ex} \hrule\vspace{2ex}

% ----------------------------------------------------- %
% LEARNING OUTCOMES
\subsection*{Learning Outcomes}
\begin{outline}
    \1      \textbf{Problem solving}\newline Identify, formulate, analyse and solve complex engineering problems in electronic engineering, creatively and innovatively.

    \1      \textbf{Application of scientific and engineering knowledge}\newline Apply knowledge of mathematics, natural sciences, engineering fundamentals and an engineering speciality to solve complex engineering problems.

    \1      \textbf{Engineering design}\newline Perform creative, procedural and non-procedural design and synthesis of components, systems, engineering works.

    \1      \textbf{Investigations, experiments and data analysis}\newline Demonstrate competence to design and conduct investigations and experiments.

    \1      \textbf{Engineering methods, skills and tools, including information technology}\newline Demonstrate competence to use appropriate engineering methods, skills and tools, including those based on information technology.

    \1      \textbf{Professional and technical communication}\newline Demonstrate competence to communicate effectively, both orally and in writing, with engineering audiences and the community at large.

    \1      \textbf{Sustainability and impact of engineering activity}\newline Demonstrate critical awareness of the sustainability and impact of engineering activity on the social, industrial and physical environment.

    \1      \textbf{Individual, team and multidisciplinary working}\newline Demonstrate competence to work effectively as an individual, in teams and in multidisciplinary environments.

    \1      \textbf{Independent learning ability}\newline Demonstrate competence to engage in independent learning through well-developed learning skills.

    \1      \textbf{Engineering professionalism}\newline Demonstrate critical awareness of the need to act professionally and ethically and to exercise judgment and take responsibility within own limits of competence.

    \1      \textbf{Engineering management}\newline Demonstrate knowledge and understanding of engineering management principles and economic decision making.

    \1      To be entrepreneurial and create job in the related industry.

    \1       Communicate well and clear, and show a good understanding of the discipline.

\end{outline}
\vspace{2ex}\hrule\vspace{2ex}

% ----------------------------------------------------- %
% Students Assenssment
\subsection*{Programme Assessment}
\begin{outline}
    \1 \textbf{Coursework} \newline
    Unless otherwise specified, continuous assessment shall contribute \qty{25}{\percent} of the final marks.\label{label1}
    \1 \textbf{Written Examination} \newline
    Unless otherwise specified, the formal examination shall contribute \qty{75}{\percent} of the final marks.
\end{outline}
\vspace{2ex}\hrule\vspace{2ex}

% ----------------------------------------------------- %
% PROGRAMME COURSES
\subsection*{Programme Summary}
\vspace{0ex}%
\noindent\begin{tcolorbox}[breakable, width=0.5\textwidth, align=left, enhanced, left= -3pt, boxrule=0pt, colback=white, colframe=white, before=\noindent, top=-13pt, bottom=-8pt]
             \noindent\begin{longtable}{@{}llS[table-format = 4.3]}
                          \textsf{\textbf{Module Code}} & \textsf{\textbf{Module Title}} & \textsf{\textbf{Credits}} \\
                          & \textbf{Part One} & \\
                          PLC1101 & Peace, Leadership and Conflict Transformation \Romannum{1}  & 5  \\
                          SCS1101 & Introduction to Computer Science and Programming & 10 \\
                          SMA1116 & Engineering Mathematics 1A & 10 \\
                          SPH1104 & Modern Physics  & 10 \\
                          TCE1103 & Professional Engineering Skills \Romannum{1} & 5 \\
                          TEE1103 & Electrical Engineering - Basic Circuit Analysis & 10\\
                          TEE1105 & Electronic Engineering Workshop & 5 \\
                          TIE1101 & Engineering Drawing \Romannum{1} & 10 \\
                          PLC1201 & Peace, Leadership and Conflict Transformation \Romannum{2} & 5\\
                          SCS1205 & Software Engineering Concepts & 10 \\
                          SMA1216 & Engineering Mathematics 1B & 10 \\
                          TEE1201 & Graphics for Electronic Engineers & 10 \\
                          TEE1203 & Electronic Engineering - Circuits and Devices & 10 \\
                          TEE1205 & Electronics Engineering Workshop & 5\\
                          TEE1211 & Analogue Electronics & 10 \\
                          TEE1241 & Electrical Measurements & 10 \\
                          & & \\
                          & \textbf{Part Two} & \\
                          SMA2116 & Engineering Mathematics \Romannum{2} & 10 \\
                          TEE2101 & Network Theory & 10 \\
                          TEE2102 & Electrical Machines & 10 \\
                          TEE2103 & Design and Project & 5 \\
                          TEE2104 & Laboratory & 5 \\
                          TEE2111 & Digital Electronics & 10 \\
                          TEE2112 & Microprocessors & 10 \\
                          SMA2217 & Engineering Mathematics \Romannum{3} & 10 \\
                          TEE2201 & Electromagnetic Fields and Materials & 10 \\
                          TEE2202 & Electronic Drives & 10 \\
                          TEE2203 & Design and Project & 5 \\
                          TEE2204 & Laboratory & 5 \\
                          TEE2205 & The Professional Engineer \Romannum{1} & 10 \\
                          TEE2211 & Digital Devices and Systems & 10 \\
                          TEE2232 & Computer Engineering \Romannum{1} & 10 \\
                          & & \\
                          & \textbf{Part Three} & \\
                          SMA3116 & Engineering Mathematics \Romannum{4} & 10 \\
                          TEE3101 & Digital Signal Processing & 10 \\
                          TEE3103 & Design and Project & 5 \\
                          TEE3104 & Laboratory & 5 \\
                          TEE3111 & Linear Integrated Circuits & 10 \\
                          TEE3121 & Analogue Communication Engineering & 10 \\
                          TEE3132 & Software Engineering & 10 \\
                          TEE3201 & Electromagnetic Theory & 10 \\
                          TEE3203 & Design and Project & 5 \\
                          TEE3204 & Laboratory & 5 \\
                          TEE3221 & Digital Communication Engineering & 10 \\
                          TEE3232 & Embedded Systems & 10 \\
                          TEE3241 & Control Engineering & 10  \\
                          & & \\
                          & \textbf{Part Four} & \\
                          TEE4000 & Industrial Attachment & 120 \\
                          & & \\
                          & \textbf{Part Five} & \\
                          TEE5101 & Engineering Management \Romannum{1} & 10 \\
                          TEE5111 & Integrated Circuit and Microelectronics & 10\\
                          TEE5133 & Communication Systems Performance & 10 \\
                          TEE5141 & Modern Control Engineering & 10 \\
                          TEE5003 & Honours Project & 50 \\
                          TEE5205 & Engineering Management \Romannum{2} & 10 \\
                          TEE5224 & High Speed Networks & 10\\
                          TEE5241 & Industrial Control & 10
             \end{longtable}%}
\end{tcolorbox}
\vspace{2ex}\hrule\vspace{2ex}
% ----------------------------------------------------- %
% PROGRAMME ACTIVITIES TIME ALLOCATION
\subsection*{Programme Activities Time Allocation}
\vspace{0ex}%
\noindent\begin{tcolorbox}[breakable, width=0.5\textwidth, align=left, enhanced, left= -3pt, boxrule=0pt, colback=white, colframe=white, before=\noindent, top=-13pt, bottom=-8pt]
             \noindent\begin{longtable}{@{}l S[table-format = 5.0] S[table-format = 5.0]}
                          \textsf{\textbf{Programme Activity}} & {\textsf{\textbf{Hours}}} & {\textsf{\textbf{Credits}}} \\
                          \multicolumn{3}{c}{\textsf{\textbf{Contact Hours}}} \\
                          Lectures & 1621 & 162 \\
                          Tutorials & 336 & 34 \\
                          Laboratory Work & 450 & 45 \\
                          Workshops & 96 & 10 \\
                          Industrial Attachment & {1 Year} & 120 \\
                          \multicolumn{3}{c}{\textsf{\textbf{Assessment Hours}}} \\
                          Final written examinations & 132 & 13 \\
                          In-class tests & 379 & 38 \\
                          Practicals & 421 & 42 \\
                          \multicolumn{3}{c}{\textsf{\textbf{Independent Study Hours}}}\\
                          Preparation for Assessments & 403 & 40 \\
                          Reading & 458 & 46 \\
                          Written assignments & 326 & 33 \\
                          Revisions & 369 & 37 \\
                          \makecell{\textsf{\textbf{Maximum Credits for \qty{80}{\percent} Modules Threshold}} } & & 620 \\
             \end{longtable}%}
\end{tcolorbox}
\vspace{2ex}\hrule
\vspace{2ex}\hrule\vspace{2ex}


